\chapter{\uppercase{System Architecture and Design}}
% Page numbering handled by Chapter 1

% ====================================================================
% CHAPTER 3: METHODOLOGY (compressed)
% ====================================================================

The literature review in Chapter~2 exposed four research gaps: (1)~static analysis tools produce false-positive rates exceeding 90\% without semantic context~\cite{ghostsecurity2025}, (2)~graph neural networks for vulnerability detection plateau at 63--65\% accuracy on single-representation inputs~\cite{devign2019}, (3)~large language models hallucinate incorrect security verdicts without grounding evidence~\cite{primevul2024}, and (4)~fixed fusion weights in hybrid systems ignore CWE-specific behavior, applying uniform combination strategies across vulnerability categories. Each gap corresponds to a strength of a different analysis modality. This chapter specifies the four modules of the proposed framework, the uncertainty-driven routing logic between them, and the formal guarantees that bind the stages together.

% ====================================================================
\section{\uppercase{System Overview}}
\label{sec:system-overview}
% ====================================================================

The framework instantiates a four-stage cascade where each finding enters at the cheapest stage and escalates only when the current stage cannot resolve it with sufficient confidence. Stage~1 (\textbf{SAST Pre-Screening}) applies Tree-sitter pattern matching and CodeQL taint analysis; each finding receives a four-factor uncertainty score $U$, and those with $U < 0.5$ resolve immediately. Stage~2 (\textbf{Graph-Augmented Validation}) builds a Joern Code Property Graph, backward-slices it, produces GraphCodeBERT embeddings concatenated with six structural features, and classifies via MiniGINv3; a conformal prediction layer wraps the classifier so that singleton prediction sets resolve while ambiguous sets escalate. Stage~3 (\textbf{LLM Dual-Agent Validation}) dispatches an Attacker agent and a Defender agent via Gemini~2.5, fuses their verdicts through a five-rule consensus protocol, and computes CVSS~v3.1 scores. Stage~4 (\textbf{Score Fusion and Reporting}) applies CWE-adaptive weighted fusion and emits SARIF~2.1.0 output with an interactive HTML dashboard. Table~\ref{tab:stages} summarizes the cost and resolution profile of each stage; Figure~\ref{fig:full-arch} provides a detailed view of data flow between all four modules.

\begin{table}[H]
\centering
\small
\caption{Stage comparison: latency, cost, and expected resolution rate.}
\label{tab:stages}
\begin{tabular}{llllr}
\toprule
\textbf{Stage} & \textbf{Engine} & \textbf{Latency} & \textbf{Cost} & \textbf{Resolution} \\
\midrule
1: SAST         & Tree-sitter + CodeQL & 100~ms + 5~min & CPU only      & 85\% \\
2: Graph        & Joern + MiniGINv3    & 2--10~s/finding & CPU/GPU       & 2\% \\
3: LLM          & LLM Provider     & 5--15~s/finding  & API tokens    & 13\% \\
4: Fusion       & Weighted aggregation & $<$1~ms         & CPU only      & 100\% \\
\bottomrule
\end{tabular}
\end{table}

\begin{figure}[H]
\centering
\resizebox{\textwidth}{!}{%
\begin{tikzpicture}[
    node distance=0.5cm and 0.7cm,
    every node/.style={font=\footnotesize},
    stage/.style={draw, rounded corners=4pt, minimum height=1.1cm, minimum width=3cm, align=center, thick},
    s1/.style={stage, fill=green!12},
    s2/.style={stage, fill=orange!12},
    s3/.style={stage, fill=red!10},
    s4/.style={stage, fill=purple!12},
    decision/.style={draw, diamond, aspect=2, fill=yellow!15, inner sep=1pt, font=\scriptsize, align=center},
    arr/.style={-Stealth, thick, color=gray!70!black},
]
\node[s1] (sast) {\textbf{Stage 1: SAST}\\{\scriptsize Tree-sitter + CodeQL}\\{\scriptsize Uncertainty $U$}};
\node[decision, right=0.7cm of sast] (d1) {$U \geq 0.5$?};
\node[s2, right=0.7cm of d1] (graph) {\textbf{Stage 2: Graph}\\{\scriptsize MiniGINv3}\\{\scriptsize Conformal}};
\node[decision, right=0.7cm of graph] (d2) {Ambiguous?};
\node[s3, right=0.7cm of d2] (llm) {\textbf{Stage 3: LLM}\\{\scriptsize Dual-Agent}\\{\scriptsize Consensus}};
\node[s4, right=0.7cm of llm] (fusion) {\textbf{Stage 4: Fusion}\\{\scriptsize CWE-adaptive}\\{\scriptsize SARIF}};
\node[left=0.5cm of sast] (input) {Source Code};
\node[below=0.6cm of d1, font=\scriptsize, text=green!40!black] (r1) {Resolved (85\%)};
\node[below=0.6cm of d2, font=\scriptsize, text=green!40!black] (r2) {Resolved (2\%)};
\draw[arr] (input) -- (sast);
\draw[arr] (sast) -- (d1);
\draw[arr] (d1) -- node[above, font=\scriptsize]{Yes} (graph);
\draw[-Stealth, thick, dashed, color=green!50!black] (d1) -- (r1);
\draw[arr] (graph) -- (d2);
\draw[arr] (d2) -- node[above, font=\scriptsize]{Yes} (llm);
\draw[-Stealth, thick, dashed, color=green!50!black] (d2) -- (r2);
\draw[arr] (llm) -- (fusion);
\node[above=0.08cm of sast, font=\scriptsize, text=gray] {$<$100\,ms};
\node[above=0.08cm of graph, font=\scriptsize, text=gray] {2--10\,s};
\node[above=0.08cm of llm, font=\scriptsize, text=gray] {5--15\,s};
\end{tikzpicture}%
}
\caption{Four-stage cascade architecture with uncertainty-driven escalation. Findings enter at Stage~1 and escalate only when the current stage cannot resolve them. Dashed arrows indicate resolution; solid arrows indicate escalation.}
\label{fig:full-arch}
\end{figure}

% ====================================================================
\section{\uppercase{Stage 1: SAST Pre-Screening}}
\label{sec:sast}
% ====================================================================

The first stage generates initial candidate findings and computes a principled uncertainty score that determines whether each finding can be resolved immediately or must escalate (Contribution~C1).

\subsection{Pre-Screening and Taint Analysis}

Before invoking CodeQL (database construction takes 1--5 minutes), a Tree-sitter pre-screening pass completes in under 100~ms per file. The scanner parses source into a concrete syntax tree and matches 24 vulnerability-indicative patterns across five languages (Python, JavaScript, Java, C/C++, Go), producing findings with baseline confidence of 0.6--0.8. When CodeQL databases are available, the \texttt{security-extended} query suite traces data flow from sources through transformations to sinks, producing findings with full taint-flow paths (confidence 0.8--0.95). Each \texttt{TaintFlow} object records ordered \texttt{TaintFlowStep} entries with \texttt{Location} and kind labels (\texttt{source}, \texttt{intermediate}, \texttt{sanitizer}, \texttt{sink}). The taint-flow length and interprocedural depth feed directly into the uncertainty scorer.

\subsection{Four-Factor Uncertainty Scoring}

A single SAST confidence score is insufficient for routing: a finding might have high confidence yet involve a 7-hop interprocedural taint path that the tool cannot fully verify. We decompose uncertainty into four independent factors:

\begin{align}
U = w_{\text{conf}} \cdot U_{\text{conf}} + w_{\text{comp}} \cdot U_{\text{comp}} + w_{\text{nov}} \cdot U_{\text{nov}} + w_{\text{confl}} \cdot U_{\text{confl}} + \delta_{\text{sev}}
\label{eq:uncertainty}
\end{align}

\noindent where $w_{\text{conf}} = 0.4$, $w_{\text{comp}} = 0.3$, $w_{\text{nov}} = 0.2$, $w_{\text{confl}} = 0.1$, and $\delta_{\text{sev}}$ is a severity adjustment. The result is clamped to $[0, 1]$. The four factors are: (1)~$U_{\text{conf}} = 1 - c_{\text{sast}}$, the inverse of SAST confidence; (2)~$U_{\text{comp}} = (\text{hop\_score} + \text{depth\_score})/2$, where $\text{hop\_score} = \text{clamp}((l_{\text{taint}} - 1)/4)$ and $\text{depth\_score} = \text{clamp}(d_{\text{interproc}}/5)$; (3)~$U_{\text{nov}} \in \{0.15, 0.85, 1.0\}$, a discrete CWE rarity classification (25 common CWEs receive 0.15, rare CWEs receive 0.85, absent CWE receives 1.0); and (4)~$U_{\text{confl}} \in \{0.0, 0.1, 0.5, 1.0\}$, measuring inter-tool disagreement.

\begin{table}[H]
\centering
\small
\caption{Uncertainty factors with weights, value ranges, and severity adjustments.}
\label{tab:uncertainty-factors}
\begin{tabular}{llcl}
\toprule
\textbf{Factor} & \textbf{Symbol} & \textbf{Weight} & \textbf{Range / Value} \\
\midrule
Confidence  & $U_{\text{conf}}$  & 0.4 & $[0, 1]$: inverse of SAST confidence \\
Complexity  & $U_{\text{comp}}$  & 0.3 & $[0, 1]$: avg of hop and depth scores \\
Novelty     & $U_{\text{nov}}$   & 0.2 & $\{0.15, 0.85, 1.0\}$: CWE rarity \\
Conflict    & $U_{\text{confl}}$ & 0.1 & $\{0.0, 0.1, 0.5, 1.0\}$: tool agreement \\
\midrule
\multicolumn{4}{l}{\textbf{Severity Adjustment} ($\delta_{\text{sev}}$, applied post-hoc to the weighted sum)} \\
\midrule
\multicolumn{2}{l}{Critical} & $+0.15$ & Force escalation for high-impact findings \\
\multicolumn{2}{l}{High}     & $+0.10$ & Lean toward deeper analysis \\
\multicolumn{2}{l}{Medium}   & $\phantom{+}0.00$ & Neutral baseline \\
\multicolumn{2}{l}{Low}      & $-0.05$ & Reduce noise from low-impact findings \\
\bottomrule
\end{tabular}
\end{table}

Pseudocode for all three algorithms (uncertainty scoring, APS calibration/inference, and consensus protocol) appears in the Appendix.

\subsection{Escalation Routing}

The \texttt{EscalationRouter} applies four rules; a finding escalates if \emph{any} rule fires.

\begin{table}[H]
\centering
\small
\caption{SAST-to-Graph escalation rules. A finding escalates if any rule fires.}
\label{tab:routing-rules}
\begin{tabular}{clp{7cm}}
\toprule
\textbf{Rule} & \textbf{Condition} & \textbf{Rationale} \\
\midrule
R1 & $U \geq 0.5$ & High composite uncertainty; SAST alone is unreliable \\
R2 & $l_{\text{taint}} > 3$ & Long taint paths exceed SAST verification capacity \\
R3 & \texttt{is\_interprocedural} = True & Cross-file flows require structural analysis \\
R4 & Interprocedural AND \texttt{CRITICAL} & Critical cross-file findings always escalate \\
\bottomrule
\end{tabular}
\end{table}

Findings that do not trigger any rule resolve at Stage~1 with verdict \texttt{SAFE}. In evaluation across 15 repositories, 85\% of findings resolved at Stage~1. For the Graph-to-LLM transition, the router applies a single rule: if the conformal prediction set is ambiguous (contains both ``safe'' and ``vulnerable''), the finding escalates; singleton sets resolve at Stage~2.

% ====================================================================
\section{\uppercase{Stage 2: Graph-Augmented Validation}}
\label{sec:graph}
% ====================================================================

Stage~2 addresses a key SAST limitation: taint paths lack structural context. A missing bounds check, an uncalled sanitizer, or a dead-code branch requires graph-level analysis. This section covers CPG construction, backward slicing, feature extraction, GNN classification, and the conformal prediction layer (Contribution~C2).

\subsection{CPG Construction and Backward Slicing}

We use Joern~\cite{cpg2014} to generate Code Property Graphs that unify 14 node types and 6 edge types (Abstract Syntax Tree edges, Control Flow Graph edges, Data Dependency Graph (DDG) edges, Control Dependency Graph (CDG) edges, CALL edges, and REACHING\_DEF edges) across Python, JavaScript, Java, C, C++, and Go. When Joern is unavailable, the builder returns a minimal stub graph (graceful degradation). Full CPGs contain thousands of irrelevant nodes, so we apply backward slicing~\cite{cpg2014} from the finding's sink via breadth-first search (BFS) along DDG, REACHING\_DEF, and CDG edges up to depth~10, then expand by $\pm 5$ lines of AST context. Typical graph reduction is 67--91\%: a 2,000-node CPG is sliced to 180--660 nodes.

\subsection{Feature Extraction}

Each node in the sliced CPG receives a 774-dimensional feature vector: 768 dimensions from GraphCodeBERT~\cite{graphcodebert2021} (\texttt{microsoft/graphcodebert-base}, CLS token, batches of 32, LRU cache of 4,096) and 6 structural features.

\begin{table}[H]
\centering
\small
\caption{Six structural node features concatenated with GraphCodeBERT embeddings.}
\label{tab:structural-features}
\begin{tabular}{llll}
\toprule
\textbf{Feature} & \textbf{Range} & \textbf{Computation} & \textbf{Security Signal} \\
\midrule
\texttt{in\_degree\_norm}  & $[0, 1]$ & $\text{in\_deg}(v) / \max(\text{in\_deg})$ & Data convergence point \\
\texttt{out\_degree\_norm} & $[0, 1]$ & $\text{out\_deg}(v) / \max(\text{out\_deg})$ & Data fan-out point \\
\texttt{is\_sink}          & $\{0, 1\}$ & Match against 54 sink patterns & Dangerous operation \\
\texttt{is\_source}        & $\{0, 1\}$ & Match against 42 source patterns & Untrusted input \\
\texttt{depth\_norm}       & $[0, 1]$ & Breadth-first search depth / max depth & Call stack position \\
\texttt{language\_id}      & $\{0, \ldots, 0.8\}$ & py=0, js=0.2, java=0.4, c=0.6, go=0.8 & Language-specific patterns \\
\bottomrule
\end{tabular}
\end{table}

The final node feature vector is $\mathbf{x}_v = [\mathbf{e}_v \| \mathbf{s}_v] \in \mathbb{R}^{774}$, where $\mathbf{e}_v \in \mathbb{R}^{768}$ is the GraphCodeBERT embedding and $\mathbf{s}_v \in \mathbb{R}^{6}$ is the structural feature vector.

\subsection{MiniGINv3 Architecture}

We chose the Graph Isomorphism Network (GIN)~\cite{xu2019gin} over GAT based on expressiveness: Xu et al.\ proved that GNNs with mean or max aggregation are strictly less expressive than the Weisfeiler-Leman (1-WL) test, while GIN's sum aggregation is as powerful as 1-WL. For vulnerability detection, code graphs often differ in subtle structural ways (a missing bounds check removes one edge, dead code adds unreachable nodes), and GIN's injective sum aggregation preserves these differences.

\begin{table}[H]
\centering
\small
\caption{GAT vs.\ GIN: theoretical and practical comparison.}
\label{tab:gat-vs-gin}
\begin{tabular}{lll}
\toprule
\textbf{Property} & \textbf{GAT (V2)} & \textbf{GIN (MiniGINv3)} \\
\midrule
Aggregation        & Weighted mean (attention) & Sum (injective) \\
WL-test equivalence & Strictly weaker & As powerful as 1-WL \\
Parameters         & 298K & 2,375,046 (2.4M) \\
Input dimension    & 773 (5 structural) & 774 (6 structural) \\
Pooling            & Global mean & Dual (mean + add) \\
\bottomrule
\end{tabular}
\end{table}

\begin{table}[H]
\centering
\small
\caption{MiniGINv3 layer-by-layer architecture (2,375,046 parameters).}
\label{tab:miniginv3-layers}
\begin{tabular}{lll}
\toprule
\textbf{Layer} & \textbf{Operation} & \textbf{Output Dim} \\
\midrule
Projection  & Linear(774, 384) + BatchNorm + ReLU & $N \times 384$ \\
GIN Layer 1 & GINConv(MLP: 384$\to$768$\to$384) + BN + ReLU + Drop(0.35) + Residual & $N \times 384$ \\
GIN Layer 2 & GINConv(MLP: 384$\to$768$\to$384) + BN + ReLU + Drop(0.35) + Residual & $N \times 384$ \\
GIN Layer 3 & GINConv(MLP: 384$\to$768$\to$384) + BN + ReLU + Drop(0.35) + Residual & $N \times 384$ \\
Dual Pool   & cat(GlobalMeanPool, GlobalAddPool) & 768 \\
Classifier  & Linear(768, 384) + ReLU + Linear(384, 2) + Softmax & 2 \\
Confidence  & Linear(768, 1) + Sigmoid & 1 \\
\bottomrule
\end{tabular}
\end{table}

Each GIN layer applies the update rule:

\begin{align}
\mathbf{h}_v^{(k)} = \text{MLP}^{(k)}\!\left((1 + \epsilon^{(k)}) \cdot \mathbf{h}_v^{(k-1)} + \sum_{u \in \mathcal{N}(v)} \mathbf{h}_u^{(k-1)}\right)
\label{eq:gin-update}
\end{align}

\noindent where $\epsilon^{(k)}$ is learnable (\texttt{train\_eps=True}), $\mathcal{N}(v)$ is the neighbor set, and $\text{MLP}^{(k)}$ is a two-layer perceptron (384$\to$768$\to$384) with BatchNorm and ReLU. Residual connections add the input after each layer. Dual pooling concatenates global mean and global add pooling into a 768-dimensional graph embedding: mean pooling captures average node behavior while add pooling preserves total feature magnitude.

The confidence head is trained with an auxiliary loss: $\mathcal{L}_{\text{conf}} = 0.1 \times \text{BCE}(\hat{c}, \mathbb{1}[\arg\max(\hat{y}) = y])$, training it to estimate prediction reliability independently of softmax output. The total loss is:
\[
\mathcal{L}_{\text{total}} = \mathcal{L}_{\text{CE}}(\hat{y}, y; w_{\text{vuln}}=1.5) + 0.1 \times \mathcal{L}_{\text{conf}}
\]

\subsection{GNN Training Configuration}

Training uses AdamW (lr $3 \times 10^{-4}$, weight decay $10^{-3}$, gradient clip 0.5) with 5-epoch linear warmup followed by cosine decay, batch size 64, dropout 0.35, and early stopping at patience 20. The training data combines Juliet Test Suite~\cite{juliet2017}, CVEfixes~\cite{cvefixes2021}, and DiverseVul~\cite{diversevul2023}, split 60/15/15/10 for training, validation, calibration, and test (calibration reserved exclusively for conformal prediction). Graphs are capped at 300 nodes; samples require $\geq$20 characters and $\geq$3 lines.

The deployed V5 model differs from V3 in two ways: dropout reduced from 0.4 to 0.35, and label smoothing removed entirely (from 0.1 in V3). Label smoothing was removed because it shifted the logit distribution, increasing nonconformity scores and producing wider prediction sets.

\subsection{Conformal Prediction Layer}

Softmax outputs are poorly calibrated~\cite{guocalibration2017}. We implement Adaptive Prediction Sets (APS)~\cite{romano2020} with $\alpha = 0.1$ (90\% coverage guarantee) to convert GNN outputs into prediction sets with formal guarantees~\cite{angelopoulos2023, vovk2005}.

\textbf{Conformal Temperature Scaling.} Before computing softmax probabilities, we apply temperature scaling with $T = 0.95$ (deployed value, tuned via grid search in Chapter~4):

\begin{align}
\hat{\pi}_i = \frac{\exp(z_i / T)}{\sum_j \exp(z_j / T)}
\label{eq:confts}
\end{align}

\textbf{Calibration.} On a held-out calibration set, we compute nonconformity scores: sort classes by descending softmax probability, compute cumulative sums $C_j = \sum_{k=1}^{j} \hat{\pi}_{(k)}$, and set $s_i = C_{\text{rank}(y_i)}$. The quantile threshold is:

\begin{align}
\hat{q} = \text{Quantile}\!\left(\{s_1, \ldots, s_n\},\; \frac{\lceil (n+1)(1-\alpha) \rceil}{n}\right)
\label{eq:aps-quantile}
\end{align}

\textbf{Coverage Guarantee.} For any new test sample $(X, Y)$ exchangeable with the calibration data:

\begin{align}
P(Y \in \mathcal{C}(X)) \geq 1 - \alpha
\label{eq:coverage}
\end{align}

This guarantee is distribution-free and holds regardless of model accuracy~\cite{barber2023}. The only assumption is exchangeability between calibration and test samples. At inference, we compute temperature-scaled probabilities, sort classes in descending order, and include classes greedily until the cumulative sum meets $\hat{q}$. Singleton sets (\{``safe''\} or \{``vulnerable''\}) resolve at Stage~2; ambiguous sets (\{``safe'', ``vulnerable''\}) escalate to Stage~3.

% ====================================================================
\section{\uppercase{Stage 3: LLM Dual-Agent Validation}}
\label{sec:llm}
% ====================================================================

Stage~3 applies LLM reasoning to the hardest cases. A single LLM judge produces unreliable security verdicts (47--76\% inconsistency across repeated tasks~\cite{steenhoek2024}), but an adversarial protocol with opposing agents achieves higher accuracy through structured disagreement~\cite{multiagentdebate2023}. This section describes the dual-agent architecture, consensus protocol (Contribution~C3), CVSS scoring, and RAG knowledge base.

\begin{figure}[H]
\centering
\resizebox{\textwidth}{!}{%
\begin{tikzpicture}[
    node distance=0.8cm and 1.0cm,
    every node/.style={font=\footnotesize},
    box/.style={draw, rounded corners=3pt, minimum height=0.85cm, align=center, thick},
    input/.style={box, fill=green!10, text width=2cm},
    db/.style={draw, rounded corners=2pt, fill=green!8, minimum height=0.55cm, align=center, font=\scriptsize, text width=1.5cm},
    proc/.style={box, fill=orange!12, text width=2.6cm},
    agent/.style={box, fill=red!10, text width=2.8cm},
    defend/.style={box, fill=blue!10, text width=2.8cm},
    engine/.style={box, fill=yellow!15, text width=2.6cm},
    output/.style={box, fill=purple!12, text width=2.4cm},
    arr/.style={-{Stealth[length=2.5mm]}, thick, color=gray!70!black},
    lbl/.style={font=\scriptsize, fill=white, inner sep=1pt},
]
% Row 1: Input
\node[input] (finding) {Escalated\\Finding};

% Row 1: RAG + databases
\node[proc, right=0.8cm of finding] (rag) {RAG Retrieval\\{\scriptsize FAISS + BM25}};
\node[db, above=0.35cm of rag] (nvd) {NVD (200K+)};
\node[db, below=0.35cm of rag] (cwe) {CWE (900+)};

% Row 1: Prompt
\node[proc, right=0.8cm of rag] (prompt) {Prompt Assembly\\{\scriptsize CWE template}};

% Row 2: Agents (stacked vertically, aligned right of prompt)
\node[agent, right=1.0cm of prompt, yshift=1.0cm] (atk) {\textbf{Attacker}\\{\scriptsize Exploit paths}};
\node[defend, right=1.0cm of prompt, yshift=-1.0cm] (def) {\textbf{Defender}\\{\scriptsize Sanitizers}};

% LLM label
\node[font=\scriptsize, fill=gray!8, rounded corners=2pt, draw=gray!40, inner sep=2pt, right=1.0cm of prompt] (llm) {LLM};

% Row 3: Consensus
\node[engine, right=1.0cm of llm, xshift=1.2cm] (consensus) {\textbf{Consensus}\\{\scriptsize 5 Rules}};

% Row 3: CVSS
\node[proc, right=0.8cm of consensus] (cvss) {CVSS v3.1};

% Row 3: Output
\node[output, right=0.8cm of cvss] (out) {Verdict +\\CVSS Score};

% Arrows
\draw[arr] (finding) -- (rag);
\draw[arr] (nvd) -- (rag);
\draw[arr] (cwe) -- (rag);
\draw[arr] (rag) -- (prompt);
\draw[arr] (prompt) -| (atk.west);
\draw[arr] (prompt) -| (def.west);
\draw[arr, color=red!50!black] (atk.east) -| ([xshift=-0.3cm]consensus.north);
\draw[arr, color=blue!50!black] (def.east) -| ([xshift=0.3cm]consensus.south);
\draw[arr] (consensus) -- (cvss);
\draw[arr] (cvss) -- (out);

% Labels
\node[lbl, text=red!60!black, above right=-0.05cm and 0.1cm of atk.east] {\textit{\scriptsize AttackerVerdict}};
\node[lbl, text=blue!60!black, below right=-0.05cm and 0.1cm of def.east] {\textit{\scriptsize DefenderVerdict}};
\end{tikzpicture}%
}
\caption{Stage~3 LLM architecture: escalated findings receive RAG context, then undergo parallel adversarial analysis by Attacker and Defender agents. A five-rule consensus engine reconciles their verdicts and computes CVSS~v3.1 base scores.}
\label{fig:llm-arch}
\end{figure}

\subsection{Dual-Agent Architecture}

The \texttt{ConsensusEngine} creates two specialized agents. The \textbf{Attacker Agent} (red team) attempts to construct a working exploit, outputting an \texttt{AttackerVerdict} with exploitability flag, concrete payload, CVSS exploitability sub-metrics (AV, AC, PR, UI), blocking factors, and confidence. The \textbf{Defender Agent} (blue team) searches for existing protections, outputting a \texttt{DefenderVerdict} with sanitizers found, access controls, framework protections, path feasibility, CVSS impact sub-metrics (S, C, I, A), and defense coverage score $d_{\text{cov}} \in [0, 1]$.

Both agents use the LLM provider (15 requests per minute, 500 requests per day on the free tier) with JSON-mode output and temperature 0.1. The adversarial framing follows multi-agent debate~\cite{divergentthinking2023}: structurally opposed agents produce more factually grounded outputs than a single ``thorough'' agent.

\subsection{CWE-Specific Prompts and Tiers}

Jinja2 templates are organized by CWE category, selected via a CWE-to-category mapping. Table~\ref{tab:cwe-templates} shows the six categories.

\begin{table}[H]
\centering
\small
\caption{CWE-to-template category mapping for Jinja2 prompt selection.}
\label{tab:cwe-templates}
\begin{tabular}{llp{5.5cm}}
\toprule
\textbf{Category} & \textbf{CWE IDs} & \textbf{Prompt Focus} \\
\midrule
\texttt{injection}       & 74, 77, 78, 79, 89, 90, 94, 95, 917 & Source-sink taint, sanitizer bypass \\
\texttt{deserialization}  & 502 & Gadget chains, type confusion \\
\texttt{path\_traversal} & 22, 23, 35, 59 & Canonicalization, chroot escape \\
\texttt{crypto}          & 327, 328, 330, 331 & Algorithm strength, key management \\
\texttt{auth}            & 284, 285, 287, 862, 863 & Session handling, privilege escalation \\
\texttt{default}         & All others & General vulnerability analysis \\
\bottomrule
\end{tabular}
\end{table}

Prompt sizing is uncertainty-adaptive through three tiers, as shown in Table~\ref{tab:prompt-tiers}.

\begin{table}[H]
\centering
\small
\caption{Prompt tiers based on uncertainty score.}
\label{tab:prompt-tiers}
\begin{tabular}{llrl}
\toprule
\textbf{Tier} & \textbf{Condition} & \textbf{Tokens} & \textbf{Contents} \\
\midrule
Minimal  & $U < 0.3$ & $\sim$500 & Code snippet + CWE name \\
Standard & $0.3 \leq U < 0.6$ & $\sim$1,500 & Code + taint path + CWE description \\
Full     & $U \geq 0.6$ & $\sim$3,000 & Everything + RAG context + graph metrics \\
\bottomrule
\end{tabular}
\end{table}

\subsection{Consensus Protocol}

The consensus engine applies five decision rules. Let $e$ denote the attacker's exploitable flag, $c_{\text{atk}}$ the attacker's confidence, $d_{\text{cov}}$ the defender's defense coverage, and $p_{\text{feas}}$ the defender's path feasibility flag.

\begin{table}[H]
\centering
\small
\caption{Consensus decision rules with configurable thresholds.}
\label{tab:consensus-rules}
\begin{tabular}{clll}
\toprule
\textbf{Rule} & \textbf{Condition} & \textbf{Verdict} & \textbf{Confidence} \\
\midrule
R1  & $e = \text{True} \wedge d_{\text{cov}} < 0.5$ & CONFIRMED & $\max(c_{\text{atk}}, 1 - d_{\text{cov}})$ \\
R2  & $e = \text{False} \wedge d_{\text{cov}} > 0.7$ & SAFE & $\max(d_{\text{cov}}, 1 - c_{\text{atk}})$ \\
R2b & $e = \text{False} \wedge p_{\text{feas}} = \text{False}$ & SAFE & $0.8$ (fixed) \\
R3  & $e = \text{True} \wedge d_{\text{cov}} \geq 0.5$ & LIKELY & $0.5 + 0.3(c_{\text{atk}} - d_{\text{cov}})$, clamped $[0.3, 0.85]$ \\
R4  & $e = \text{False} \wedge d_{\text{cov}} \leq 0.7$ & POTENTIAL & $0.4 + 0.2(1 - d_{\text{cov}})$, clamped $[0.2, 0.6]$ \\
\bottomrule
\end{tabular}
\end{table}

\subsection{CVSS v3.1 Scoring}

The framework computes CVSS~v3.1 base scores from eight sub-metrics produced by the two agents, following the NIST specification~\cite{cvss31}. The attacker provides four exploitability sub-metrics: Attack Vector (AV), Attack Complexity (AC), Privileges Required (PR), and User Interaction (UI). The defender provides Scope (S) and three impact sub-metrics: Confidentiality (C), Integrity (I), and Availability (A). The Impact Sub-Score $\text{ISS} = 1 - (1-C)(1-I)(1-A)$ feeds into scope-dependent impact formulas, while exploitability is $8.22 \times AV \times AC \times PR \times UI$:

\begin{align}
\text{ISS} &= 1 - (1 - C)(1 - I)(1 - A) \label{eq:iss} \\
\text{Exploitability} &= 8.22 \times AV \times AC \times PR \times UI \label{eq:exploitability} \\
\text{Impact}_{\text{unchg}} &= 6.42 \times \text{ISS} \\
\text{Impact}_{\text{chg}} &= 7.52 \times (\text{ISS} - 0.029) - 3.25 \times (\text{ISS} - 0.02)^{15} \\
\text{Base Score} &= \begin{cases}
    \text{roundup}(\min(\text{Impact} + \text{Exploitability},\; 10.0)) & \text{if } S = \text{unchanged} \\
    \text{roundup}(\min(1.08 \times (\text{Impact} + \text{Exploitability}),\; 10.0)) & \text{if } S = \text{changed}
\end{cases} \label{eq:cvss-base}
\end{align}

\noindent For findings resolved before Stage~3, default CVSS vectors are provided for 14 common CWEs. Table~\ref{tab:cvss-vectors} shows representative defaults.

\begin{table}[H]
\centering
\small
\caption{Default CVSS v3.1 vectors for common CWEs. Legend: N=Network/None, L=Low/Local, H=High, R=Required, U=Unchanged, C=Changed.}
\label{tab:cvss-vectors}
\begin{tabular}{lccccccccr}
\toprule
\textbf{CWE} & \textbf{AV} & \textbf{AC} & \textbf{PR} & \textbf{UI} & \textbf{S} & \textbf{C} & \textbf{I} & \textbf{A} & \textbf{Score} \\
\midrule
CWE-78 (OS Cmd Inj.)  & Net & Low & None & None & Unchg & High & High & High & 9.8 \\
CWE-89 (SQL Injection) & Net & Low & None & None & Unchg & High & High & None & 9.1 \\
CWE-79 (XSS)           & Net & Low & None & Req  & Chgd  & Low  & Low  & None & 6.1 \\
CWE-416 (Use After Free)& Loc & High& Low  & None & Unchg & High & High & High & 7.0 \\
CWE-327 (Broken Crypto) & Net & Low & None & None & Unchg & High & None & None & 7.5 \\
CWE-502 (Deserialization)& Net& Low & None & None & Unchg & High & High & High & 9.8 \\
\bottomrule
\end{tabular}
\end{table}

\subsection{RAG Knowledge Base}

Table~\ref{tab:rag-config} summarizes the knowledge base configuration. The knowledge base indexes over 200,000 NVD CVE entries and over 900 CWE entries. Retrieval combines Facebook AI Similarity Search (FAISS) semantic search (\texttt{all-MiniLM-L6-v2}, flat inner-product index) with BM25 keyword search, fused via Reciprocal Rank Fusion (Equation~\ref{eq:rrf}). The top-5 documents after RRF fusion are injected into agent prompts; batch processing groups up to 5 findings per LLM call, reducing API calls by up to 5$\times$.

\begin{align}
\text{RRF}(d) = \sum_{r \in \{r_{\text{FAISS}}, r_{\text{BM25}}\}} \frac{w_r}{k + \text{rank}_r(d)}
\label{eq:rrf}
\end{align}

\begin{table}[H]
\centering
\small
\caption{RAG knowledge base configuration.}
\label{tab:rag-config}
\begin{tabular}{ll}
\toprule
\textbf{Parameter} & \textbf{Value} \\
\midrule
Semantic model & \texttt{all-MiniLM-L6-v2} \\
FAISS index type & \texttt{IndexFlatIP} (inner product) \\
BM25 variant & Okapi BM25 \\
Semantic weight ($w_{\text{FAISS}}$) & 0.6 \\
Keyword weight ($w_{\text{BM25}}$) & 0.4 \\
RRF constant ($k$) & 60 \\
Top-$k$ results & 5 \\
NVD entries & 200,000+ \\
CWE entries & 900+ \\
OWASP mappings & 40+ CWEs $\to$ OWASP Top 10 (2021) \\
\bottomrule
\end{tabular}
\end{table}

% ====================================================================
\section{\uppercase{Stage 4: Score Fusion and Classification}}
\label{sec:fusion}
% ====================================================================

Stage~4 combines confidence scores from all stages into a unified vulnerability score using CWE-adaptive weighting (Contribution~C4).

\subsection{CWE-Adaptive Fusion Formula}

\begin{align}
S_{\text{fused}} = \alpha \cdot S_{\text{SAST}} + \beta \cdot S_{\text{Graph}} + \gamma \cdot S_{\text{LLM}}
\label{eq:fusion}
\end{align}

\noindent where $(\alpha, \beta, \gamma)$ are CWE-specific weights from \texttt{configs/cwe\_weights.yaml}. When a finding did not pass through a given stage, the engine renormalizes remaining weights to sum to~1.0. Table~\ref{tab:cwe-weights} lists the configured weights.

\begin{table}[H]
\centering
\small
\caption{Representative CWE-adaptive fusion weights ($\alpha$, $\beta$, $\gamma$). The complete set of 14 profiles appears in the Appendix, Table~\ref{tab:fusion-weights}.}
\label{tab:cwe-weights}
\begin{tabular}{llcccp{4cm}}
\toprule
\textbf{Category} & \textbf{CWE IDs} & $\alpha$ & $\beta$ & $\gamma$ & \textbf{Rationale} \\
\midrule
Injection      & 78, 79        & 0.25 & 0.25 & 0.50 & LLM excels at context \\
SQL Injection  & 89            & 0.30 & 0.25 & 0.45 & SAST taint strong here \\
Code Injection & 94            & 0.25 & 0.30 & 0.45 & Graph structure matters \\
Crypto         & 327, 328      & 0.50 & 0.20 & 0.30 & SAST pattern matching \\
Memory         & 416           & 0.20 & 0.50 & 0.30 & Graph captures structure \\
Null Deref     & 476           & 0.25 & 0.45 & 0.30 & Graph captures structure \\
Auth           & 287, 862      & 0.20 & 0.25 & 0.55 & LLM semantic understanding \\
Path Traversal & 22            & 0.35 & 0.30 & 0.35 & Balanced taint + context \\
Deserialization & 502          & 0.25 & 0.25 & 0.50 & LLM understands gadgets \\
Default        & All others    & 0.30 & 0.30 & 0.40 & Balanced baseline \\
\bottomrule
\end{tabular}
\end{table}

Two LLM override rules supersede threshold classification: (1)~if the attacker confirmed exploitability and defense coverage $< 0.3$, the verdict is forced to CONFIRMED; (2)~if the attacker found no exploit and defense coverage $> 0.8$, the verdict is forced to SAFE.

\subsection{Three-Tier Classification}

\begin{table}[H]
\centering
\small
\caption{Classification thresholds applied to the fused score.}
\label{tab:classification}
\begin{tabular}{lcl}
\toprule
\textbf{Verdict} & \textbf{Threshold} & \textbf{Interpretation} \\
\midrule
CONFIRMED  & $S_{\text{fused}} \geq 0.85$ & High-confidence true positive \\
LIKELY     & $0.50 \leq S_{\text{fused}} < 0.85$ & Probable vulnerability \\
POTENTIAL  & $0.0 < S_{\text{fused}} < 0.50$ & Low confidence, needs review \\
SAFE       & (via LLM override) & False positive filtered \\
\bottomrule
\end{tabular}
\end{table}

% ====================================================================
\section{\uppercase{Finding Lifecycle and Reporting}}
\label{sec:reporting}
% ====================================================================

Each finding follows a defined lifecycle encoded in the \texttt{Finding} Pydantic model, which serves as the universal data container across all four stages. Every intermediate value is preserved, making the entire analysis chain auditable in the final report.

Findings transition through seven states: NEW, SAST\_SCORED, SAST\_RESOLVED, GRAPH\_VALIDATED, GRAPH\_RESOLVED, LLM\_VALIDATED, and SCORE\_FUSED. A finding enters as NEW when created by Tree-sitter or CodeQL, receives an \texttt{UncertaintyScore} at SAST\_SCORED, and resolves at the earliest stage whose confidence is sufficient. Findings that reach LLM\_VALIDATED carry attacker/defender verdicts, a consensus result, and a CVSS score before proceeding to final fusion.

The framework emits SARIF~2.1.0 output with custom properties recording uncertainty scores, stage resolution, fused scores, and CVSS vectors. The \texttt{--dashboard} flag generates a self-contained HTML report with a cascade funnel, expandable finding cards, severity distribution, CWE treemap, and CVSS histogram. The pipeline degrades gracefully when tools are missing: CodeQL-unavailable uses Tree-sitter alone, Joern-unavailable skips Stage~2, missing LLM keys skip Stage~3.

\vspace{0.5em}

This chapter defined the complete pipeline from source code to triaged vulnerability report. Chapter~4 evaluates the pipeline experimentally: cascade efficiency, classification accuracy, conformal coverage, and end-to-end comparison against standalone tools.
